\chapter{Continuous Random Variables}\label{ch:g12:contdist}

Waiting times, physical measurements, proportions: many random quantities
take a continuum of values, and no single value has positive \omterm{def:g10:proba:distribution}{probability}.
Probabilities are then computed by integrating a \emph{\omterm{def:g12:contdist:density}{density}}. This
chapter studies the uniform, exponential and \omterm{def:g12:contdist:normal}{normal distributions}, and uses
the normal law to quantify the fluctuations of polls and \omterm{def:g10:proba:sample}{samples}.

\section{Probability densities}

\begin{definition}[Density, continuous random variable]\label{def:g12:contdist:density}
A \emph{probability density}\index{density} on an \omterm{def:g10:numbers:interval}{interval} $I$ is a
\omterm{def:g12:limcont:continuity}{continuous}, nonnegative \omterm{def:g11:func:function}{function} $f$ on $I$ with $\int_I f(t)\,\dd t = 1$
(the \omterm{def:g12:integ:area}{integral} over an unbounded $I$ being understood as a limit of
\omterm{def:g12:integ:area}{integrals} over growing bounded \omterm{def:g10:numbers:interval}{intervals}). A \omterm{def:g12:randvar:rv}{random variable} $X$ has
density $f$ if for all $a \leq b$ in $I$:
\[
\P(a \leq X \leq b) = \int_a^b f(t)\,\dd t .
\]
\end{definition}

Probabilities are \emph{areas under the \omterm{def:g12:contdist:density}{density} curve}. In particular
$\P(X = a) = \int_a^a f = 0$ for every single value $a$: only \omterm{def:g10:numbers:interval}{intervals}
carry \omterm{def:g10:proba:distribution}{probability}, and $\P(a \leq X \leq b) = \P(a < X < b)$.

\begin{definition}[Expectation and variance]\label{def:g12:contdist:exp}
For $X$ with \omterm{def:g12:contdist:density}{density} $f$ on $I$:
\[
\E(X) = \int_I t\,f(t)\,\dd t, \qquad
\V(X) = \int_I \bigl(t - \E(X)\bigr)^2 f(t)\,\dd t
= \E(X^2) - \E(X)^2 .
\]
\end{definition}

The rules of Chapters 14--15 (linearity of \omterm{def:g12:randvar:exp}{expectation}, additivity of
\omterm{def:g12:randvar:exp}{variance} for \omterm{def:g12:sums:indep}{independent} variables, Bienaymé--Chebyshev) remain valid;
their proofs, with \omterm{def:g12:integ:area}{integrals} in place of sums, are admitted at this level.

\section{Uniform distribution}

\begin{definition}[Uniform distribution]\label{def:g12:contdist:uniform}
$X$ follows the \emph{uniform distribution}\index{uniform distribution}
$\mathcal U\intcc ab$ if its \omterm{def:g12:contdist:density}{density} is constant,
$f(t) = \frac{1}{b-a}$ on $\intcc{a}{b}$. Then for
$\intcc{c}{d} \subseteq \intcc ab$,
$\P(c \leq X \leq d) = \frac{d - c}{b - a}$: \omterm{def:g10:proba:distribution}{probability} proportional to
length.
\end{definition}

\begin{proposition}\label{prop:g12:contdist:uniformmoments}
If $X \sim \mathcal U\intcc ab$, then
$\E(X) = \dfrac{a+b}{2}$ and $\V(X) = \dfrac{(b-a)^2}{12}$.
\end{proposition}

\begin{proof}
$\E(X) = \int_a^b \frac{t}{b-a}\dd t
= \frac{b^2 - a^2}{2(b-a)} = \frac{a+b}{2}$. For the \omterm{def:g12:randvar:exp}{variance},
$\E(X^2) = \frac{b^3 - a^3}{3(b-a)} = \frac{a^2 + ab + b^2}{3}$, and
\[
\V(X) = \frac{a^2 + ab + b^2}{3} - \frac{(a+b)^2}{4}
= \frac{4a^2 + 4ab + 4b^2 - 3a^2 - 6ab - 3b^2}{12}
= \frac{(b - a)^2}{12}. \qedhere
\]
\end{proof}

\section{Exponential distribution}

\begin{definition}[Exponential distribution]\label{def:g12:contdist:expo}
For $\lambda > 0$, $X$ follows the \emph{exponential
distribution}\index{exponential distribution} $\mathcal E(\lambda)$ if its
\omterm{def:g12:contdist:density}{density} on $\intco{0}{+\infty}$ is
\[
f(t) = \lambda\,\eu^{-\lambda t}.
\]
Then $\P(X \leq x) = 1 - \eu^{-\lambda x}$ and
$\P(X > x) = \eu^{-\lambda x}$ for $x \geq 0$.
\end{definition}

\begin{proposition}\label{prop:g12:contdist:expomoments}
If $X \sim \mathcal E(\lambda)$:
$\E(X) = \dfrac{1}{\lambda}$ and $\V(X) = \dfrac{1}{\lambda^2}$.
\end{proposition}

\begin{proof}
Integrating by parts on $\intcc{0}{A}$:
\[
\int_0^A t\,\lambda\eu^{-\lambda t}\dd t
= \bigl[-t\,\eu^{-\lambda t}\bigr]_0^A + \int_0^A \eu^{-\lambda t}\dd t
= -A\eu^{-\lambda A} + \frac{1 - \eu^{-\lambda A}}{\lambda}
\xrightarrow[A\to+\infty]{} \frac1\lambda,
\]
using $A\eu^{-\lambda A} \to 0$ (\cref{thm:g12:exp:growth}). A second
\omterm{thm:g12:integ:ibp}{integration by parts} gives $\E(X^2) = \frac{2}{\lambda^2}$, whence
$\V(X) = \frac{2}{\lambda^2} - \frac{1}{\lambda^2}$.
\end{proof}

\begin{theorem}[Memorylessness]\label{thm:g12:contdist:memoryless}
\index{memorylessness}
If $X \sim \mathcal E(\lambda)$, then for all $s, t \geq 0$:
\[
\pcond{X > s}{X > s + t} = \P(X > t) .
\]
The \omterm{def:g12:contdist:expo}{exponential distribution} is the law of lifetimes \emph{without aging}
(radioactive nuclei, not light bulbs).
\end{theorem}

\begin{proof}
\[
\pcond{X > s}{X > s+t}
= \frac{\P(X > s + t)}{\P(X > s)}
= \frac{\eu^{-\lambda(s+t)}}{\eu^{-\lambda s}} = \eu^{-\lambda t}
= \P(X > t). \qedhere
\]
\end{proof}

\begin{omfigure}
\begin{tikzpicture}
\begin{axis}[omaxis,
  width=8.8cm, height=5cm,
  xmin=-0.4, xmax=4.4, ymin=0, ymax=1.15,
  xtick={1.3}, xticklabels={$t$}, ytick={1}, yticklabels={$\lambda$}]
  \addplot[name path=dens, omDef, thick, domain=0:4.2, samples=80]
    {exp(-x)};
  \path[name path=xaxis] (0,0) -- (4.2,0);
  \addplot[omThm!20] fill between[of=dens and xaxis,
    soft clip={domain=1.3:4.2}];
  \draw[black, dashed, thin] (axis cs:1.3,0) -- (axis cs:1.3,0.2725);
  \node[omThm, font=\small, anchor=west] at (axis cs:2.1,0.35)
    {$\P(X > t) = \eu^{-\lambda t}$};
\end{axis}
\end{tikzpicture}

{\small The exponential \omterm{def:g12:contdist:density}{density} $\lambda\eu^{-\lambda x}$ (here
$\lambda = 1$): the tail area beyond $t$ (red) is $\eu^{-\lambda t}$,
and memorylessness says every tail looks like the whole \omterm{def:g12:randvar:rv}{distribution}
rescaled.}
\end{omfigure}

\section{Normal distribution}

\begin{definition}[Normal distribution]\label{def:g12:contdist:normal}
$X$ follows the \emph{standard normal distribution}\index{normal
distribution} $\mathcal N(0, 1)$ if its \omterm{def:g12:contdist:density}{density} on $\R$ is
\[
\varphi(t) = \frac{1}{\sqrt{2\pi}}\,\eu^{-t^2/2}
\]
(the bell curve). More generally,
$X \sim \mathcal N(\mu, \sigma^2)$ if $\dfrac{X - \mu}{\sigma} \sim
\mathcal N(0,1)$; then $\E(X) = \mu$ and $\V(X) = \sigma^2$.
\end{definition}

\begin{omfigure}
\begin{tikzpicture}
\begin{axis}[
  width=11cm, height=4.6cm,
  axis lines=middle, xmin=-4, xmax=4, ymin=0, ymax=0.45,
  xtick={-3,-2,-1,0,1,2,3},
  xticklabels={$\mu{-}3\sigma$,$\mu{-}2\sigma$,$\mu{-}\sigma$,$\mu$,
               $\mu{+}\sigma$,$\mu{+}2\sigma$,$\mu{+}3\sigma$},
  ytick=\empty, samples=120, domain=-4:4]
  \addplot[name path=gauss, omDef, thick] {exp(-x^2/2)/sqrt(2*pi)};
  \path[name path=xaxis] (-4,0) -- (4,0);
  \addplot[omDef!20] fill between[of=gauss and xaxis,
           soft clip={domain=-2:2}];
  \node[omDef] at (axis cs:0,0.17) {$\approx 95.4\%$};
\end{axis}
\end{tikzpicture}

{\small The normal \omterm{def:g12:contdist:density}{density}: about $68.3\%$ of the mass lies within
$\sigma$ of the \omterm{def:g11:stat:mean}{mean}, $95.4\%$ within $2\sigma$, $99.7\%$ within
$3\sigma$.}
\end{omfigure}

\begin{remark}
The factor $\frac{1}{\sqrt{2\pi}}$ makes the total area $1$ --- a famous
computation (the \emph{Gaussian \omterm{def:g12:integ:area}{integral}}) done at university. There is no
elementary formula for $\int \varphi$: normal probabilities are read from
tables or calculators.
\end{remark}

\begin{theorem}[De Moivre--Laplace, admitted]\label{thm:g12:contdist:tcl}
\index{De Moivre--Laplace theorem}
Let $X_n \sim \mathcal B(n, p)$ and
$Z_n = \dfrac{X_n - np}{\sqrt{np(1-p)}}$ (the standardized binomial). Then
for all $a \leq b$,
\[
\P(a \leq Z_n \leq b) \xrightarrow[n\to+\infty]{}
\int_a^b \varphi(t)\,\dd t .
\]
\emph{This result is admitted at this level.} It explains the universal
appearance of the bell curve: a binomial with large $n$ is approximately
normal --- and (central limit theorem, university) so is any sum of many small
\omterm{def:g12:sums:indep}{independent} effects.
\end{theorem}

\begin{method}[Fluctuation interval at $95\%$]\label{met:g12:contdist:fluctuation}
\index{confidence interval}
Since a normal variable falls within $1.96$ \omterm{def:g11:stat:variance}{standard deviations} of its \omterm{def:g11:stat:mean}{mean}
with \omterm{def:g10:proba:distribution}{probability} $0.95$, for $n$ large a \omterm{def:g10:stats:series}{frequency}
$F_n = \frac{X_n}{n}$ of a $\mathcal B(n,p)$ \omterm{def:g10:proba:sample}{sample} satisfies
\[
\P\left(p - 1.96\sqrt{\tfrac{p(1-p)}{n}} \leq F_n \leq
p + 1.96\sqrt{\tfrac{p(1-p)}{n}}\right) \approx 0.95 .
\]
\begin{itemize}
  \item \emph{Fluctuation \omterm{def:g10:numbers:interval}{interval}} (testing): if a claimed $p$ puts the
        observed \omterm{def:g10:stats:series}{frequency} outside this \omterm{def:g10:numbers:interval}{interval}, reject the claim at the
        $5\%$ level.
  \item \emph{\omterm{met:g12:contdist:fluctuation}{Confidence interval}} (estimating): the simplified \omterm{def:g10:numbers:interval}{interval}
        $\left[F_n - \frac{1}{\sqrt n},\ F_n + \frac{1}{\sqrt n}\right]$
        contains $p$ with \omterm{def:g10:proba:distribution}{probability} at least $0.95$
        (using $1.96\sqrt{p(1-p)} \leq 1$, see \cref{exo:g12:sums:8}).
\end{itemize}
\end{method}

\begin{example}\label{ex:g12:contdist:poll}
A poll of $n = 1000$ voters gives a candidate $52\%$. The \omterm{met:g12:contdist:fluctuation}{confidence
interval} $52\% \pm \frac{1}{\sqrt{1000}} \approx 52\% \pm 3.2\%$ still
contains $50\%$: the poll alone does not establish that the candidate is
ahead.
\end{example}

\section{Exercises}

\begin{exercise}[$\star$]\label{exo:g12:contdist:1}
A bus passes every $15$ minutes; a traveler arrives at a uniformly random
time. Let $X \sim \mathcal U\intcc{0}{15}$ be the waiting time. Compute
$\P(X \leq 5)$, $\P(4 \leq X \leq 10)$, $\E(X)$ and $\sigma(X)$.
\end{exercise}

\begin{exercise}[$\star$]\label{exo:g12:contdist:2}
Check that $f(t) = \frac{3}{8}t^2$ is a \omterm{def:g12:contdist:density}{density} on $\intcc{0}{2}$, and
compute $\E(X)$ and $\P(X \geq 1)$ for a variable $X$ with this \omterm{def:g12:contdist:density}{density}.
\end{exercise}

\begin{exercise}[$\star$]\label{exo:g12:contdist:3}
The lifetime (in years) of an electronic component follows
$\mathcal E(0.2)$.
\begin{enumerate}
  \item Compute the expected lifetime and $\P(X > 5)$.
  \item The component has already worked $3$ years. What is the
        \omterm{def:g10:proba:distribution}{probability} it works at least $5$ more?
\end{enumerate}
\end{exercise}

\begin{exercise}[$\star\star$]\label{exo:g12:contdist:4}
The half-life of a radioactive atom whose lifetime follows
$\mathcal E(\lambda)$ is the \omterm{def:g11:stat:median}{median} $m$: $\P(X > m) = \frac12$. Express
$m$ as a \omterm{def:g11:func:function}{function} of $\lambda$ and compare with the \omterm{def:g12:randvar:exp}{expectation}. Which is
larger, and why does that make sense for a skewed \omterm{def:g12:randvar:rv}{distribution}?
\end{exercise}

\begin{exercise}[$\star\star$]\label{exo:g12:contdist:5}
Heights in a population follow $\mathcal N(175,\ 7^2)$ (in cm). Using the
$68$--$95$--$99.7$ rule, estimate the proportion of the population with
height between $168$ and $182$ cm, above $189$ cm, and below $154$ cm.
\end{exercise}

\begin{exercise}[$\star\star$]\label{exo:g12:contdist:6}
A machine fills bags labeled $500$ g; the mass filled follows
$\mathcal N(\mu,\ 4^2)$. Regulations demand that at most $2.5\%$ of bags
weigh less than $500$ g. Using the $95\%$ rule, what minimal setting of
$\mu$ complies?
\end{exercise}

\begin{exercise}[$\star\star$]\label{exo:g12:contdist:7}
A die claimed to be fair is rolled $1200$ times and shows a six $260$
times (as in \cref{exo:g12:sums:7}).
\begin{enumerate}
  \item Compute the $95\%$ fluctuation \omterm{def:g10:numbers:interval}{interval} for the \omterm{def:g10:stats:series}{frequency} of sixes
        of a fair die over $1200$ rolls.
  \item Is the observed \omterm{def:g10:stats:series}{frequency} inside? Compare the strength of this
        conclusion with the Bienaymé--Chebyshev analysis.
\end{enumerate}
\end{exercise}

\begin{exercise}[$\star\star\star$]\label{exo:g12:contdist:8}
Before an election, a poll of $n$ people will estimate the score $p$ of a
candidate by the observed \omterm{def:g10:stats:series}{frequency} $F_n$.
\begin{enumerate}
  \item With the \omterm{met:g12:contdist:fluctuation}{confidence interval} of
        \cref{met:g12:contdist:fluctuation}, what \omterm{def:g10:proba:sample}{sample} size guarantees a
        margin of $\pm 2$ points?
  \item The candidates are separated by $1$ point in reality. Explain why
        no realistic poll can reliably call the winner, however well
        conducted.
\end{enumerate}
\end{exercise}

\begin{exercise}[$\star\star\star$]\label{exo:g12:contdist:9}
Let $X$ have \omterm{def:g12:contdist:density}{density} $f$ on $\intcc ab$ and let
$Y = \alpha X + \beta$ with $\alpha > 0$.
\begin{enumerate}
  \item Show that $\P(c \leq Y \leq d)
        = \P\left(\frac{c - \beta}{\alpha} \leq X \leq
        \frac{d-\beta}{\alpha}\right)$, and deduce that $Y$ has \omterm{def:g12:contdist:density}{density}
        $g(y) = \frac{1}{\alpha} f\left(\frac{y - \beta}{\alpha}\right)$.
  \item Deduce that if $Z \sim \mathcal N(0,1)$, then
        $\mu + \sigma Z$ has \omterm{def:g12:contdist:density}{density}
        $\frac{1}{\sigma\sqrt{2\pi}}\,
        \eu^{-(y - \mu)^2/(2\sigma^2)}$ --- the general normal \omterm{def:g12:contdist:density}{density}.
\end{enumerate}
\end{exercise}
